\begin{abstract}
Rolling element bearings are critical components in industrial machinery, and their failure can lead to catastrophic system breakdowns. While deep learning has significantly advanced fault diagnosis, the performance of these models is often hampered by the class imbalance problem, where fault samples are scarce compared to normal data. To address this challenge, this paper proposes a novel framework termed CAC-CycleGAN-WGP for imbalanced bearing fault diagnosis. The methodological core involves a CycleGAN architecture that transforms normal vibration signals into realistic fault signals, integrated with a Conditional Auxiliary Classifier (CAC) to accurately capture the marginal probability distributions of various fault categories. To ensure training stability and suppress mode collapse, we incorporate Wasserstein distance with a Gradient Penalty (WGP). Furthermore, a quality assessment method based on Stacked Autoencoders (SAE) is developed to evaluate the fidelity of the generated samples. Extensive experiments are conducted on two imbalanced bearing datasets: the Case Western Reserve University (CWRU) dataset and a custom laboratory dataset. Comparative results against state-of-the-art methods, including SMOTE, GAN, WGAN, ACGAN, and CAC-CycleGAN-WGP, demonstrate that the proposed CAC-CycleGAN-WGP significantly improves diagnostic accuracy and robustness under severe data imbalance, providing a viable solution for real-world industrial health monitoring.
\end{abstract}

\section{Introduction}
\label{sec:intro}

Mechanical equipment in modern industrial systems is becoming increasingly complex and automated, where rolling element bearings serve as fundamental yet vulnerable components. The health status of bearings directly affects the reliability and safety of the entire production line. Consequently, intelligent fault diagnosis has emerged as a pivotal research area in the field of Prognostics and Health Management (PHM). Recent years have witnessed the transformative impact of deep learning on diagnostic performance, as highlighted in several comprehensive surveys \cite{8, Lei2020, 9}. Various architectures, including convolutional neural networks (CNNs), recurrent neural networks (RNNs), and hybrid CNN-LSTM models \cite{10}, have demonstrated strong feature extraction capabilities for vibration signal analysis. These methods leverage deep architectures to automatically extract representative features from raw vibration signals, surpassing traditional signal processing techniques in many scenarios \cite{32}. However, the effectiveness of these data-driven models is predicated on the availability of balanced and annotated datasets, a requirement that is rarely met in practical industrial environments \cite{40, 41}. To ensure the safety of large-scale manufacturing, it is essential to develop systems that can recognize impending failures even when historical fault data is extremely limited.

Despite the proliferation of monitoring sensors, the data collected from industrial machinery is inherently imbalanced. Machinery typically operates in a normal state for the vast majority of its service life, making fault samples extremely rare and difficult to acquire. This data scarcity leads to a significant class imbalance problem, which poses a severe challenge to standard deep learning classifiers. In real-world industrial settings, the ratio of fault samples to normal samples often ranges from 1:10 to 1:100, creating an extremely skewed distribution that hinders model convergence \cite{32, 40}. The rarity of fault instances means that the decision boundaries of standard classifiers are heavily biased towards the healthy state, potentially leading to missed detections of critical failures. For instance, in wind turbine maintenance or high-speed rail monitoring, a single misclassified bearing defect can result in millions of dollars in downtime or risk to human lives. This data scarcity leads to a significant class imbalance problem, which poses a severe challenge to standard deep learning classifiers. When trained on such skewed distributions, models tend to overfit the majority (normal) class and fail to generalize to the minority (fault) classes, resulting in poor diagnostic sensitivity for the very failures they are designed to detect. The robustness of fault diagnosis systems under these imbalanced conditions remains a critical bottleneck for their deployment in safety-critical applications.

To mitigate the effects of class imbalance, various data augmentation and resampling techniques have been proposed. Traditional methods such as the Synthetic Minority Over-sampling Technique (SMOTE) \cite{6} generate synthetic samples by interpolating between existing minority instances. Extensions of this approach, such as feature-level SMOTE \cite{18}, attempt to address this limitation by performing oversampling in a learned feature space rather than the raw signal domain. While SMOTE and its derivatives have been widely adopted \cite{1, 2, 5, 35}, they often struggle with high-dimensional temporal data like vibration signals. These methods can introduce noise and fail to capture the complex underlying physics of bearing dynamics, leading to synthetic samples that lack the necessary diversity or realistic characteristics required for training sophisticated deep neural networks. Furthermore, traditional statistical oversampling assumes a local linearity that does not hold for the impulsive and non-stationary nature of mechanical vibration signals, often resulting in "overlapping" features between different fault types which further confuses the diagnostic classifier. Recently, researchers have turned to deep generative models to learn the actual manifold of the fault signals, aiming to create more physically meaningful augmentations.

The advent of Generative Adversarial Networks (GANs) \cite{1} offered a more powerful alternative for generating high-fidelity synthetic data. By employing a minimax game between a generator and a discriminator, GANs can learn the distribution of fault signals more effectively than traditional interpolation. Various GAN architectures have been explored for PHM, including Deep Convolutional GAN (DCGAN) \cite{14} which uses spatial filters to capture local correlations, and Conditional GAN (CGAN) which introduces class labels to guide the generation process. Several recent works have demonstrated the effectiveness of GAN-based augmentation for imbalanced bearing fault diagnosis \cite{33, 43}. Additionally, WGAN uses the Earth Mover distance to provide a more stable optimization objective, while attention-enhanced variants of conditional GANs \cite{34} have further improved generation quality under severe imbalance. However, vanilla GANs and their variants, such as WGAN \cite{2} and ACGAN \cite{5}, are prone to several well-documented issues, including mode collapse and training instability. Although Wasserstein GAN with Gradient Penalty (WGAN-GP) \cite{3} was introduced to stabilize the training process by enforcing a Lipschitz constraint on the discriminator, existing GAN-based fault diagnosis methods \cite{36, 37} often lack the ability to perform cross-domain transformation—specifically, the ability to leverage the abundant information in normal signals to synthesize rare fault patterns. CycleGAN-based methods have recently emerged as a promising direction for this cross-domain translation problem in fault diagnosis \cite{39}. Furthermore, traditional GANs often fail to maintain class-specific consistency during the generation process \cite{4, 36}.

In this paper, we propose a novel generative framework, CAC-CycleGAN-WGP, specifically designed to address imbalanced bearing fault diagnosis through signal-to-signal translation. The architecture is illustrated in Fig. 1. Unlike traditional GANs that generate samples from random noise, our approach utilizes a CycleGAN \cite{4} structure to transform readily available normal signals into synthetic fault signals, thereby preserving the structural information inherent in the vibration data. The technical motivation for using WGP is to provide stable gradients even when the support of the real and generated distributions is disjoint, which is common in high-dimensional signal spaces. By using the Wasserstein objective, we ensure that the generator receives a continuous learning signal even when the synthetic samples are initially far from the target domain. By incorporating the CAC module, we explicitly align the generator's output with the diagnostic labels, ensuring that the synthesized signals are not only "realistic" in a statistical sense but also "diagnostically accurate" for specific fault categories like inner-race or outer-race defects. This prevents the "label-flipping" problem where a generator might produce a signal that looks real but represents the wrong fault type. To enhance the discriminative power and class-specificity of the generated samples, we integrate a Conditional Auxiliary Classifier (CAC), which allows the model to learn and utilize the marginal probability distributions of different fault categories. To ensure the reliability of the synthesis, we adopt the Wasserstein distance coupled with a Gradient Penalty (WGP) to avoid mode collapse and ensure stable convergence. Additionally, we introduce an evaluation metric based on Stacked Autoencoders (SAE) \cite{44} to quantitatively assess the quality and diversity of the generated signals, ensuring they are suitable for augmenting the training set of a downstream classifier.

The primary contributions of this work are summarized as follows:
\begin{enumerate}
    \item We propose the CAC-CycleGAN-WGP method, which combines cross-domain signal translation with conditional auxiliary classification to generate high-quality fault samples from normal signals.
    \item We integrate Wasserstein distance and gradient penalty into the CycleGAN framework, overcoming the limitations of training instability and mode collapse common in previous generative fault diagnosis models.
    \item We develop a Stacked Autoencoder (SAE) based assessment strategy to objectively evaluate the fidelity of generated synthetic vibration signals.
    \item We validate the effectiveness of our approach on two distinct datasets: the CWRU benchmark (Case I) and a realistic laboratory dataset (Case II). The results demonstrate superior performance compared to SMOTE, GAN, WGAN, ACGAN, CAC-CycleGAN-WGP \cite{35}, and other competitive methods such as ICAC-CycleGAN-WGP \cite{26} and ACGAN-SDAE \cite{21}.
\end{enumerate}

The remainder of this paper is organized as follows. Section II provides a detailed review of the related literature concerning imbalanced learning, generative adversarial networks, and their applications in mechanical health monitoring. We focus on recent advances in stable GAN training and domain adaptation techniques for sensor data. Section III describes the proposed CAC-CycleGAN-WGP methodology in detail, including the architecture design, loss functions, and the SAE-based evaluation metric. We provide a rigorous mathematical formulation of the adversarial and cycle-consistency losses. Section IV presents the experimental setup, results, and a comprehensive discussion of the proposed method's performance across different datasets and imbalance ratios. We include detailed ablation studies and sensitivity analyses to validate each component. Finally, Section V concludes the paper by summarizing the key findings and suggesting directions for future research in automated industrial fault diagnosis, focusing on real-time deployment and edge computing.
